\documentclass[a4paper,11pt]{article}
\usepackage{exam-style-341}
\examinehead{341 农业综合知识三 · 数据库技术与应用}{模拟试卷（一）参考答案}
\begin{document}
\answercover{341 农业综合知识三 · 数据库技术与应用}{模拟试卷（一）参考答案}

本答案按试卷题号顺序给出，每部分均标注分值。选择题、SQL 题与规范化题附有解析，规范化题给出候选码推导、闭包运算、分解过程与无损连接性验证。

% ============================================================
\answerpart{一、单项选择题（10 题，每题 2 分，共 20 分）}
% ============================================================

\q \textbf{B}。
\begin{analysis}
数据库系统的三级模式结构为外模式、模式、内模式。外模式（也称为子模式或用户模式）是数据库用户能够看见和使用的局部数据的逻辑结构，一个数据库可以有多个外模式；模式（也称为逻辑模式）是数据库中全体数据的全局逻辑结构，一个数据库只有一个模式；内模式（也称为存储模式）是数据物理结构和存储方式的描述，一个数据库只有一个内模式。故选 B。存储模式即内模式的别名，与选项 D 同义，但题干问的是“全体数据的全局逻辑结构”，只有模式符合。
\end{analysis}

\q \textbf{B}。
\begin{analysis}
两级映像的作用：外模式/模式映像定义外模式与模式之间的对应关系，当模式改变时，DBA 修改该映像即可使外模式保持不变，由于应用程序是依据外模式编写的，因而应用程序不必修改，从而保证了数据的逻辑独立性；模式/内模式映像定义模式与内模式之间的对应关系，当内模式改变时，修改该映像可使模式保持不变，从而保证了数据的物理独立性。故选 B。
\end{analysis}

\q \textbf{A}。
\begin{analysis}
数据模型的三要素是数据结构、数据操作和完整性约束。数据结构描述数据库的组成对象以及对象之间的联系，是对系统静态特性的描述；数据操作是对数据库中各种对象（型）的实例（值）允许执行的操作的集合，包括操作及有关的操作规则，是对系统动态特性的描述；完整性约束是一组完整性规则，用以保证数据的正确、有效和相容。选项 B 是 E-R 图（概念模型）的三个组成部分，选项 C 是数据库三级模式，选项 D 是关系模型中的基本术语。
\end{analysis}

\q \textbf{B}。
\begin{analysis}
候选码（候选键）是能够唯一标识一个元组且不含多余属性的属性组；若一个关系有多个候选码，则选定其中一个作为主码（主键）；主属性是包含在任何一个候选码中的属性；超码（超键）是能够唯一标识元组但可能含有多余属性的属性组。题干强调“不含多余属性”，故为候选码。
\end{analysis}

\q \textbf{A}。
\begin{analysis}
实体完整性规则：若属性（组）$A$ 是基本关系 $R$ 的主属性，则 $A$ 不能取空值。参照完整性规则针对外码，要求外码的值或者为空值，或者等于被参照关系中某个元组的主码值。用户定义完整性是用户针对具体应用环境定义的约束条件。域完整性指属性取值必须来自其值域。故选 A。
\end{analysis}

\q \textbf{B}。
\begin{analysis}
在 $R(A,B,C)$ 上，$A^{+}=ABC$，所以 $A$ 是唯一的候选码。非主属性 $B$、$C$ 完全函数依赖于码 $A$，满足 2NF 的要求；但 $F$ 中存在 $A\to B$、$B\to C$，于是 $A\to C$ 是经 $B$ 传递而来的，即非主属性 $C$ 传递函数依赖于码 $A$，不满足 3NF。故 $R$ 最高达到 2NF。
\end{analysis}

\q \textbf{B}。
\begin{analysis}
自然连接是一种特殊的等值连接：它要求两个关系中进行比较的分量必须是同名属性组，并且在结果中把重复的属性列去掉。$R(A,B)$ 与 $S(B,C)$ 的公共属性为 $B$，连接条件为 $R.B=S.B$，结果关系模式为 $(A,B,C)$，共 3 个属性。
\end{analysis}

\q \textbf{D}。
\begin{analysis}
ACID 的含义：原子性（Atomicity）指事务中的操作要么全做要么全不做；一致性（Consistency）指事务执行的结果必须使数据库从一个一致性状态变到另一个一致性状态；隔离性（Isolation）指一个事务的执行不能被其他事务干扰；持久性（Durability）指事务一旦提交，其对数据库中数据的改变就应该是永久性的。题干描述的是持久性。
\end{analysis}

\q \textbf{C}。
\begin{analysis}
数据库设计分为需求分析、概念结构设计、逻辑结构设计、物理结构设计、数据库实施、数据库运行与维护六个阶段。概念结构设计的任务是设计概念模型（E-R 图）；将 E-R 图转换为关系模型（关系模式）属于逻辑结构设计阶段的任务；物理结构设计确定数据的存储结构与存取方法。故选 C。
\end{analysis}

\q \textbf{A}。
\begin{analysis}
事务故障是指某个事务在运行过程中未能正常终结（如运算溢出、死锁被选中撤销等）。恢复子系统反向扫描日志文件，对尚未提交的事务执行 UNDO 操作，将其对数据库的所有修改恢复到修改前的值；REDO（重做）用于系统故障后处理已提交但修改可能仍留在缓冲区而未写回磁盘的事务；重装数据库副本是介质故障的恢复手段之一。故选 A。
\end{analysis}

% ============================================================
\answerpart{二、填空题（5 题，每题 2 分，共 10 分）}
% ============================================================

\q 第一空：\textbf{选择}；第二空：\textbf{投影}。
\begin{analysis}
五种基本关系代数运算是并（$\cup$）、差（$-$）、笛卡尔积（$\times$）、选择（$\sigma$）和投影（$\pi$）。交、连接、除等运算都可以用这五种基本运算来表达。
\end{analysis}

\q 第一空：\textbf{元组}；第二空：\textbf{属性}。
\begin{analysis}
关系模型用二维表表示关系：表中的一行即一个元组（记录），表中的一列即一个属性（字段），属性的取值范围称为域。
\end{analysis}

\q \textbf{参照}。
\begin{analysis}
关系数据库的三类完整性约束包括实体完整性、参照完整性和用户定义完整性。参照完整性规则的内容是：若属性（组）$F$ 是基本关系 $R$ 的外码，它与基本关系 $S$ 的主码 $K$ 相对应，则 $R$ 中每个元组在 $F$ 上的值必须为：或者取空值，或者等于 $S$ 中某个元组的主码值。题干描述的正是参照完整性。
\end{analysis}

\q 第一空：\textbf{$B$}；第二空：\textbf{$ABCD$}。
\begin{analysis}
闭包运算：
（1）$B^{+}$：初始为 $\{B\}$。$F$ 中的三个函数依赖 $AB\to C$、$C\to D$、$D\to A$ 的左部都含有 $B$ 以外的属性（$A$ 或 $C$ 或 $D$），都不能使用，故 $B^{+}=\{B\}$。
（2）$(AB)^{+}$：初始为 $\{A,B\}$；由 $AB\to C$ 得 $\{A,B,C\}$；由 $C\to D$ 得 $\{A,B,C,D\}$；由 $D\to A$ 未产生新属性。故 $(AB)^{+}=ABCD$，即 $AB$ 能推出全部属性，$AB$ 是 $R$ 的候选码（事实上 $R$ 的候选码共有三个：$AB$、$BC$、$BD$；单个属性 $A$、$B$、$C$、$D$ 都不能唯一标识元组）。
\end{analysis}

\q 第一空：\textbf{逻辑}；第二空：\textbf{物理}。
\begin{analysis}
数据库设计通常分为需求分析、概念结构设计、逻辑结构设计、物理结构设计、数据库实施、数据库运行与维护六个阶段。其中逻辑结构设计把概念模型转换为某个 DBMS 支持的数据模型并对其进行优化；物理结构设计为逻辑数据模型选取合适的物理存储结构与存取方法。
\end{analysis}

% ============================================================
\answerpart{三、简答题（2 题，每题 5 分，共 10 分）}
% ============================================================

\q 视图及其优点：
\begin{analysis}
\textbf{（1）视图的定义（2 分）}

视图是从一个或几个基本表（或其他视图）导出的表。视图所对应的数据并不实际存储在数据库中，数据库中只存放视图的定义，因此视图是一个\textbf{虚表}；视图中的数据仍存放在导出该视图的基本表中。对视图的一切操作最终都要转换为对基本表的操作。

\textbf{（2）使用视图的优点（答出 4 点即可得 3 分）}

1）视图能够简化用户的操作。用户可以把经常使用的多表连接、投影、分组统计等查询定义为视图，以后只需对视图进行简单查询即可。

2）视图使用户能以多种角度看待同一数据。不同用户可以根据自己的需要定义不同的视图，数据共享的灵活性更高。

3）视图对重构数据库提供了一定程度的逻辑独立性。当基本表的结构发生变化时，可以通过重新定义视图使用户的外模式保持不变，从而应用程序不必修改。

4）视图能够对机密数据提供安全保护。可以有选择地让不同用户只能通过视图访问其权限内的数据，而对基本表中的其他数据保密。

5）适当利用视图可以更清晰地表达查询，把复杂的查询分解为若干步，便于理解与维护。

（说明：视图的可更新性有一定限制，由多个基本表导出或含有分组统计、表达式的视图通常不允许直接更新，这是得分点之外的补充说明。）
\end{analysis}

\q 并发操作可能带来的三类数据不一致性问题：
\begin{analysis}
\textbf{（1）丢失修改（丢失更新）。}两个事务 $T_1$ 和 $T_2$ 读入同一数据并加以修改，$T_2$ 的提交结果破坏了 $T_1$ 提交的结果，导致 $T_1$ 的修改被丢失。（约 1.5 分）

\textbf{（2）不可重复读。}事务 $T_1$ 读取某一数据后，事务 $T_2$ 对该数据执行更新操作，使 $T_1$ 无法再现前一次读取的结果。具体包括三种情形：$T_1$ 读后 $T_2$ 对其做了修改；$T_1$ 按某条件读取一组数据后 $T_2$ 删除了其中部分记录，$T_1$ 再次按相同条件读取时发现少了记录（幻影现象）；$T_1$ 读后 $T_2$ 插入了新记录。（约 1.5 分）

\textbf{（3）读“脏”数据。}事务 $T_1$ 修改某一数据并将其写回数据库，事务 $T_2$ 随后读取该数据，而 $T_1$ 由于某种原因被撤销（ROLLBACK），其修改被回滚，此时 $T_2$ 读到的数据与数据库中的实际数据不一致，是错误的数据。（约 2 分）

并发控制机制（封锁、可串行化调度等）正是为解决上述三类不一致性问题而设置的。
\end{analysis}

% ============================================================
\answerpart{四、综合题（2 题，每题 5 分，共 10 分）}
% ============================================================

\q SQL 题参考答案：

\textbf{（1）定义表 SC（1 分）}
\begin{lstlisting}[language=SQL]
CREATE TABLE SC (
    Sno   CHAR(8),
    Cno   CHAR(4),
    Grade INT,
    PRIMARY KEY (Sno, Cno),
    FOREIGN KEY (Sno) REFERENCES Student(Sno),
    FOREIGN KEY (Cno) REFERENCES Course(Cno)
);
\end{lstlisting}

\textbf{（2）查询选修“数据库”且成绩不低于 90 分的学生（2 分）}
\begin{lstlisting}[language=SQL]
SELECT S.Sno, S.Sname, SC.Grade
FROM Student S, Course C, SC
WHERE S.Sno = SC.Sno
  AND SC.Cno = C.Cno
  AND C.Cname = '数据库'
  AND SC.Grade >= 90
ORDER BY SC.Grade DESC;
\end{lstlisting}

\textbf{（3）查询平均成绩高于 80 分且至少选修 2 门课程的学生（2 分）}
\begin{lstlisting}[language=SQL]
SELECT Sno, COUNT(*) AS 选课门数, AVG(Grade) AS 平均成绩
FROM SC
GROUP BY Sno
HAVING COUNT(*) >= 2 AND AVG(Grade) > 80
ORDER BY AVG(Grade) DESC;
\end{lstlisting}

\begin{analysis}
\textbf{第（1）问：}主码由两个属性构成，必须写成表级约束 \code{PRIMARY KEY (Sno, Cno)}，不能写成两个列级主码；两个外码分别引用 \code{Student(Sno)} 和 \code{Course(Cno)}，用 \code{FOREIGN KEY ... REFERENCES ...} 定义。若把 \code{Sno}、\code{Cno} 声明为 \code{NOT NULL} 更完整（主属性不允许取空值，实体完整性）。

\textbf{第（2）问：}课程名在 \code{Course} 中，学生姓名在 \code{Student} 中，成绩在 \code{SC} 中，因此需要 \code{Student}、\code{Course}、\code{SC} 三表连接。“成绩不低于 90 分”即 \code{Grade >= 90}；“按成绩降序”即 \code{ORDER BY SC.Grade DESC}。等价的子查询写法（相关子查询 + \code{IN}/\code{EXISTS}）也可得分：先查出“数据库”的课程号，再在 \code{SC} 中筛选。

\textbf{第（3）问：}按学号分组后，\code{COUNT(*)} 是该学生的选课门数，\code{AVG(Grade)} 是其平均成绩。“至少选修 2 门”与“平均成绩高于 80 分”都是对分组结果的筛选条件，必须写在 \code{HAVING} 子句中，写在 \code{WHERE} 中会出错，因为 \code{WHERE} 在分组之前执行，不允许使用组函数。\code{AVG} 会自动忽略 \code{NULL} 值；若要求把空成绩按 0 分参加平均计算，应写 \code{AVG(COALESCE(Grade, 0))}。排序也可使用别名：\code{ORDER BY 平均成绩 DESC}。注意 \code{GROUP BY Sno} 之后 \code{SELECT} 列表中只能出现分组属性和组函数。
\end{analysis}

\q 规范化题参考答案：

\textbf{（1）候选码与主属性（1 分）}

候选码：\textbf{(订单号, 商品号)}（唯一）。

推导（闭包运算）：
\begin{itemize}[leftmargin=1.6em, itemsep=1pt]
  \item 订单号$^{+}=\{$订单号, 客户号, 客户名$\}$（由 订单号$\to$客户号、客户号$\to$客户名）；
  \item 商品号$^{+}=\{$商品号, 商品名, 单价$\}$（由 商品号$\to$商品名、商品号$\to$单价）；
  \item 其余单属性 客户号$^{+}=\{$客户号, 客户名$\}$、客户名$^{+}=\{$客户名$\}$、单价$^{+}=\{$单价$\}$、数量$^{+}=\{$数量$\}$ 都不能推出全部属性；
  \item (订单号, 商品号)$^{+}=$ 订单号, 商品号 $\cup$ $\{$客户号, 客户名$\}$ $\cup$ $\{$商品名, 单价$\}$ $\cup$ $\{$数量$\}$ $=R$ 的全部属性。
\end{itemize}
且去掉 订单号 或 商品号 后都不能推出全部属性，故 (订单号, 商品号) 是唯一的候选码。

主属性：\textbf{订单号、商品号}；非主属性：\textbf{客户号、客户名、商品名、单价、数量}。

\textbf{（2）$R$ 最高属于第几范式（2 分）}

$R$ 最高属于 \textbf{1NF}。

理由：$R$ 的每个属性都是不可再分的原子属性，故满足 1NF。$R$ 的候选码为 (订单号, 商品号)，而 $F$ 中存在
\[
\text{订单号}\to\text{客户号},\quad \text{订单号}\to\text{客户名},\quad
\text{商品号}\to\text{商品名},\quad \text{商品号}\to\text{单价},
\]
即非主属性函数依赖于候选码的真子集（订单号 或 商品号 单独即可决定部分非主属性），这是\textbf{部分函数依赖}，违反 2NF 的条件，因此 $R$ 不属于 2NF，当然也不属于 3NF。所以 $R$ 的最高范式为 1NF。

\textbf{（3）分解为 3NF（2 分）}

第一步，求 $F$ 的最小依赖集 $F_{\min}$。$F$ 中各函数依赖的右部都是单属性；各依赖左部没有多余属性；逐个检查冗余：去掉 订单号$\to$客户号 后 订单号$^{+}=\{$订单号$\}$，不能推出客户号，故不冗余；同理其余各依赖也都不能由其他依赖推出（如去掉 (订单号,商品号)$\to$数量 后 (订单号,商品号)$^{+}$ 中不含数量）。所以 $F_{\min}=F$。

第二步，按保持函数依赖的 3NF 分解算法，对每个函数依赖 $X\to A$ 取一个关系模式 $XA$，左部相同的依赖合并，得到
\[
\rho=\{R_1, R_2, R_3, R_4\},
\]
\[
R_1(\text{订单号},\text{客户号}),\quad
R_2(\text{客户号},\text{客户名}),\quad
R_3(\text{商品号},\text{商品名},\text{单价}),\quad
R_4(\text{订单号},\text{商品号},\text{数量}).
\]

第三步，检查范式级别：$R_1$ 的码为 订单号（订单号$\to$客户号）；$R_2$ 的码为 客户号；$R_3$ 的码为 商品号（商品号$\to$商品名、商品号$\to$单价，两个依赖的左部都是码）；$R_4$ 的码为 (订单号, 商品号)。各模式中非主属性都完全函数依赖于码，且不存在非主属性对码的传递函数依赖，故 $\rho$ 的每个关系模式都属于 3NF（实际上 $R_1\sim R_4$ 均属于 BCNF）。$R_4$ 中已包含 $R$ 的候选码，因此不必再补充候选码关系模式。

第四步，验证保持函数依赖：$F_{\min}$ 中每一个函数依赖都包含在某个 $R_i$ 中——
订单号$\to$客户号 在 $R_1$ 中，客户号$\to$客户名 在 $R_2$ 中，商品号$\to$商品名、商品号$\to$单价 在 $R_3$ 中，(订单号, 商品号)$\to$数量 在 $R_4$ 中，故分解保持函数依赖。

第五步，验证无损连接性（二元分解判定定理：对 $\{R_i,R_j\}$，若 $R_i\cap R_j\to(R_i-R_j)$ 或 $R_i\cap R_j\to(R_j-R_i)$ 成立，则该步连接无损）：
\begin{itemize}[leftmargin=1.6em, itemsep=1pt]
  \item $R_4\cap R_1=\{$订单号$\}$，$R_1-R_4=\{$客户号$\}$，而 订单号$\to$客户号 $\in F^{+}$，故 $R_4$ 与 $R_1$ 连接无损，得 $T_1(\text{订单号},\text{商品号},\text{数量},\text{客户号})$；
  \item $T_1\cap R_2=\{$客户号$\}$，$R_2-T_1=\{$客户名$\}$，而 客户号$\to$客户名 $\in F^{+}$，故与 $R_2$ 连接无损，得 $T_2(\text{订单号},\text{商品号},\text{数量},\text{客户号},\text{客户名})$；
  \item $T_2\cap R_3=\{$商品号$\}$，$R_3-T_2=\{$商品名, 单价$\}$，而 商品号$\to$商品名、商品号$\to$单价 均 $\in F^{+}$，故与 $R_3$ 连接无损，最终得到 $R$ 的全部属性且不丢失元组。
\end{itemize}
因此 $\rho=\{R_1,R_2,R_3,R_4\}$ 具有无损连接性，同时保持函数依赖，且每个模式都达到 3NF，满足题目要求。

\begin{analysis}
\textbf{易错点提示：}
1）候选码必须用闭包运算验证，不能凭直觉。本题中 订单号 单独不是码（推不出商品号、商品名、单价、数量），商品号 单独也不是码。
2）判断是否满足 2NF，只需检查是否存在\textbf{非主属性对候选码的真子集的函数依赖}；本题有 4 处部分函数依赖，因此只能是 1NF。
3）分解时也可以把 $R_3$ 拆成 $R_3'(\text{商品号},\text{商品名})$ 和 $R_3''(\text{商品号},\text{单价})$ 两个模式，同样满足 3NF、保持函数依赖且无损连接；两种写法都对，但把左部相同的依赖合并到一个关系模式中更简洁。
4）无损连接性必须用定理逐步验证或列表格用 Chase 算法验证，“分解后能自然连接还原”这句话本身不算证明。
\end{analysis}

\end{document}
